% ============================================================
%  第8章 异构系统下解析OpenCL
%  源文件：content/chapter8/
% ============================================================

第2章中，我们讨论过如何权衡不同的硬件架构(支持OpenCL)。这些架构支持OpenCL就是为了能够在这种模型下对现有的程序进行加速。本章中，我们会将OpenCL映射到实际硬件上，所使用的架构为一个高端CPU和一个离散GPU。虽然这里选择AMD系统来进行介绍，不过其他厂商在同种硬件架构的实现方面都差不多。


% ---- 8.1-chinese.md ----
\section{AMD FX-8350 CPU}

AMD的OpenCL实现可以运行在AMD显卡上和所有x86架构的CPU上。所有主机端代码在x86架构的CPU上执行。不过，AMD的OpenCL实现也可以将x86架构的处理器作为设备，让x86设备运行OpenCL C代码。图8.1展示了FX-8350 CPU的内部架构，该图用来描述x86架构与OpenCL实现之间的映射关系。

要在OpenCL运行时将FX-8350 CPU作为设备，需要使用clGetDeviceIDs()获取该设备句柄。然后将设备句柄作为设备对象传入clCreateContext()，clCreateCommandQueue()和clBuildProgram()。当要使用CPU作为OpenCL C的执行设备时，需要向clGetDeviceIDs()传入CL\_\allowbreak{}DEVICE\_\allowbreak{}TYPE\_\allowbreak{}CPU标识(或CL\_\allowbreak{}DEVICE\_\allowbreak{}TYPE\_\allowbreak{}ALL)。

\myGraphic{0.85}{images/chapter8/8-1.png}{图8.1 基于AMD Piledriver(打桩机)架构的FX-8350 CPU高层结构。}

CPU中OpenCL可以使用的核数有8个。如果将整个CPU作为一个独立设备，那么可以从单个队列把并行负载分散到各个CPU核上，用来分散并行负载量，这种负载均衡的方式在系统中效率最高。当然在OpenCL中也可以使用设备划分的扩展方式，将CPU划分成多个设备。

\subsection{运行时实现}

OpenCL使用CPU时，运行时在每个核上创建了一个线程(例如，线程池)，以便OpenCL内核的执行。另一个主管理线程会根据队列上的任务，将不同的任务分配给不同的线程，同时将正在执行的任务从队列上移除。任意给定的OpenCL内核可能包含成千上万个工作组(其参数以及相应内存都需要提前准备好)。

OpenCL使用栅栏的方式提供细粒度的同步。在基于CPU的传统系统中，操作系统会对线程间通讯进行管理，与操作系统交互的开销成了并行实现高效扩展的障碍。另外，如果将一个工作组分配到多个CPU核上回造成共享缓存的问题。为了缓解这个问题，OpenCL CPU运行时实现将一个工作组部署在一个CPU核(线程)上。OpenCL工作组中的工作项会按顺序执行(串行)。当前工作组中所有工作项执行完成后，将执行下一个工作组中的工作项。这样来看的话，虽然有很多工作组可以同时执行，但是工作组中的线程并不是并行的关系。图8.2中描绘了OpenCL在FX-8350 CPU上的映射关系。

\myGraphic{0.85}{images/chapter8/8-2.png}{图8.2 OpenCL映射到FX-8350 CPU上。在该场景中，FX-8350 CPU既是OpenCL主机，又是OpenCL设备。}

OpenCL中工作项在工作组内是并发，可使用栅栏进行同步。所有线程必须到达栅栏处，才能继续下面的操作。CPU作为设备时，栅栏操作等于是一个工作项终止后，另一个工作项启动；不过，这对于操作系统不现实，因为那意味着由操作系统通过线程抢占来处理(例如，中断某个线程，让另一个线程进行)。因此，当整个工作组属于一个线程时，抢占也不会有任何作用。AMD的OpenCL CPU运行时实现中，栅栏操作会使用到setjmp和longjmp指令。setjmp将会保存系统状态(保护现场)，longjmp将会读取保存了的系统状态，继续之前的执行状态[1]。不过，运行时会为了配合硬件的分支预测器和保证程序栈对齐，针对这两个函数提供两个自定义版本。图8.3展示了CPU线程在工作组中的工作流。

\myGraphic{0.85}{images/chapter8/8-3.png}{图8.3 x86架构上工作组执行的实现。}

因为工作项的执行方式是串行的，所以当在内核中使用栅栏操作时，它会先执行组内每个工作项直到栅栏，再执行栅栏之后的每个工作项，从而提供正确的栅栏语义。

使用setjmp指令时，寄存器中保存的工作项数据会在setjmp调用时备份到主存的工作项栈中。为了减少缓存争夺和未命中的情况，并提高层级缓存的利用率，栈数据存放在缓存的何处需要仔细考虑。另外，工作项栈数据是交错存储在内存上，以减少访存冲突，并且数据保存在一个较大的内存页中，以保证物理地址映射的连续性，减少CPU使用转换内存的压力。

\subsection{工作项内的向量化操作}

AMD打桩机微架构具有128位向量寄存器，可以用来处理各种版本SSE和AVX指令。OpenCL C包括一系列向量类型：float2, float4, int4等等。对于数学操作来说，有OpenCL C提供了对应的向量版本，式例如下：

\begin{codebox}[cpp]
float4 a = intput_data[location];
float4 b = a + (float4)(0.f, 1.f, 2.f, 3.f);
output_data[location] = b;
\end{codebox}

AMD打桩机微架构中，这些向量都会保存在寄存器中，并且在编译过程中会将这些操作翻译成SSE和AVX指令。这里提供了很重要的性能优化。向量的加载和存储操作，会使用底层代码进行，以提高内存操作的效率。目前，单个工作项可以使用SIMD向量：8.2节将会看到在GPU端的操作与CPU端的区别。

\subsection{局部内存}

AMD打桩机设计中并未对暂存式内存提供专用的硬件。CPU通常会为了减少访存延迟，会有很多层缓存。局部数据通常为了更加高效的访问数据，会将其映射到CPU缓存上，开发者还是享受到了来自硬件的访存加速。为了提高缓存的命中率，局部内存在每个CPU线程上都会开辟一段，并且局部内存可以被工作组中的工作项重复使用。对于串行的工作组来说(因为栅栏、数据竞争或内存冲突)，只要没有数据竞争或内存冲突，这块局部内存就不会再产生缓存未命中。局部内存另外的好处是，减少了内存频繁开辟的开销。图8.4展示了局部内存与AMD CPU缓存的映射关系。

\myGraphic{0.85}{images/chapter8/8-4.png}{图8.4 工作组(工作组0)的内存空间映射到Piledriver CPU缓存。}

虽然，在CPU上使用局部内存也会有潜在的性能收益，不过局部内存会给一些应用带来性能上的负面影响。如果内核中数据具有很好的局部性(比如，矩阵相乘的一部分)，那么用局部内存作存储实际上会进行一次不必要的拷贝，并且数据存储在L1缓存上，访存效率极高。在这种情况下，如果可用的缓存过小，则会导致性能退化，过小的缓存会造成数据争夺缓存位置，从而增加未命中率和数据拷贝的开销。

\myRefRule

\begin{myRefs}

[1] J. Gummaraju, L.Morichetti, M. Houstion, B. Snder, B.R. Gaster, B. Zheng, Twin peaks: a software platform for heterogeneous computing on general-purpose and graphics processors. in: PACT 2010: Proceedings of the Nineteenth International Conference on Parallel Architectures and Complilation Techniques. Septerber 11-15, 2010, Vienna, Austria. Association for Computing Machinery, 2010.

\end{myRefs}



% ---- 8.2-chinese.md ----
\section{AMD RADEON R9 290X GPU}

本节我们将讨论一下AMD GPU中的“波前”(在NVIDIA显卡中称为“线程束”)。这个概念是有助于我们区分OpenCL工作项和CUDA线程。不过，有时会无可避免的使用“线程”这个词来对GPU的工作方式进行描述。本书中我们将会使用“线程”来描述硬件线程。虽然，本节指定的设备为AMD Radeon R9 290X GPU，但是硬件线程与工作项映射的关系、线程调度和抢占的考量，以及内存系统的布局，不同厂商的同类产品其实很相似。

GPU上OpenCL的目标代码与CPU上的差别很大。读者应该还记得，图像处理器主要用来对3D图像进行渲染。这种设计让使得资源的优先级有所不同，与CPU架构也有很大的区别。目前GPU市场上，GPU的一些主要特性的差距越来越小，这些特性我们在第2章讨论过：

\begin{enumerate}
\item 执行宽SIMD：大量执行单元对不同的数据，执行相同的指令。
\item 大量多线程：GPU计算核支持并发大量线程。
\item 硬件暂存式内存：物理内存可由编程者控制。
\end{enumerate}

下面罗列出来的区别则更加微妙，不过其根据工作分配和交互对应用性能进行提高：

\begin{itemize}
\item 支持硬件同步：并发线程间支持细粒度交互。
\item 硬件管理任务分配：工作队列管理与负载均和由硬件控制。
\end{itemize}

硬件同步能为计算单元中的多个波前减少同步开销，并且细粒度交互的开销非常低。

GPU为任务分发提供广阔的硬件空间，因为图像处理器支持三维图像。游戏任务在显卡的流水线上交错排布，让图形任务的管理变得更加复杂。图8.5中展示了OpenCL模型在AMD Radeon R9 290X GPU映射关系，其架构包含一个命令处理器和工作组生成器(生成工作组供硬件调度)。调度器会将计算任务分配到设备的44个核上。有关计算单元的内容将在8.2.2中详细讨论。

\myGraphic{0.85}{images/chapter8/8-5.png}{图8.5 Radeon R9 290X高层示意图，标注了OpenCL执行模型和内存模型的术语。}

为了使用GPU获得高性能加速，调度需要十分高效。对于图像处理器，波前调度的开销必须很低，因为工作块可能很小(例如由几个像素构成的单个三角形)。因此，当我们需要使用GPU获取较高的性能时，我们需要：

\begin{itemize}
\item 提供大量的内核任务供硬件分发
\item 如果内核过小，考虑将内核进行合并
\end{itemize}

为每个内核提供足够工作，这样就能持续的占用工作组的流水线，所以调度器总能收到很多的任务，并且调度器通常会将给多的任务推到SIMD单元上。所以我们希望，设备每一次的吞吐可以创建大量的波前发送给GPU。

第二点是因为OpenCL设备支持队列模式。当OpenCL运行时选择执行相关设备上任务队列中的任务时，运行时会扫描队列中的任务，并选择合适大小的任务进行执行。这种任务由一个GPU工作命令缓存构成，其命令的执行放置在GPU的前端流水线上。其执行内容包括(1)构建一个队列，(2)将该队列放置在内存的某个位置上，(3)告诉设备这个队列在哪里，(4)要求设备执行这个队列。因为操作阶段这几个步骤是串行的关系，所以对于单个工作块来说，就会导致较高的延迟。另外，GPU在驱动层之下执行内核的过程中，为了能让GPU正常工作，需要多次进出内核空间进行上下文切换。CPU端如果涉及到线程上下文切换，那应用的性能肯定不会太好。所以，如果任务量太小，就会出现多个上下文切换的情况，从而导致应用性能的下降。不过，这里有个开销常量，也就是对于单个工作队列的分配，以及之后对队列上任务执行的开销。通过传递大内核的方式(或较长的内核序列)，可以避免这种开销的累积。这样做的目的，就是为了增加每个队列中所要处理的任务量。

\subsection{线程和内存系统}

CPU缓存层级结构的设计就是为了减少内存访问流的延迟，一旦内存访问流出现了重大的延迟，那么会对性能造成很大的影响。另外，因为GPU核使用线程和宽SIMD单元，GPU以牺牲延迟为代价，用线程和宽SIMD单元来最大化吞吐。图8.6展示了一个包含FX-8350 CPU和Radeon R9 290X GPU系统的内存层级结构。

这种内存系统中，GPU具有如下特征：

\begin{itemize}
\item 大量的寄存器
\item 用于管理软件的便签式内存，在AMD硬件中称为\emph{共享本地数据}(LDS, local data shares)
\item 高级别的片上多线程
\item 单个L2缓存
\item 高带宽内存
\end{itemize}

\myGraphic{0.85}{images/chapter8/8-6.png}{图8.6 独立系统中的内存带宽。}

如之前所提到的，图像任务与计算任务有很大的区别，这就导致GPU需要有自己的执行和内存模型。实际上，GPU比CPU更加少的依赖数据重用，并且GPU的缓存大小要比CPU的小很多。即使在x86核和Radeon R9 290X计算单元上的L1缓存数据大小相同，GPU一次执行40个波前的话，每个波前所分配到的缓存更加的有限。GPU对于缓存的弱依赖是由一系列原因的，其中包括时域图像数据不可复用和由于数据集过大和大量的线程，无法保证为每个工作集提供足够的缓存空间。当数据复用发生时，多线程和高带宽内存将有助于克服缓存空间有限的情况。

整个GPU只有一个多存储体的L2缓存，为所有计算单元提供数据。计算单元中的单L2缓存是为了保证GPU中L1缓存和外部缓存数据一致性而设计。这种写直达设计使得寄存器溢出变得不可取，那么后续的每次访问则会有很高的延迟，并且访存操作会阻塞在L2缓存处。

为了避免溢出，GPU提供了大量的寄存器。例如，x86架构下通用寄存器的数量是固定的(每个线程上下文中有16个通用寄存器)，Radeon R9 290X显卡单个波前可用的寄存器数量为16384个！波前尝试只使用寄存器进行计算，并且也会使用到LDS，也会尽可能避免访存冲突。

LDS允许高带宽低延迟(编程可控)的读写访问。这种可编程模式的数据复用要比硬件控制更加高效。减少垃圾数据(加载到缓存而从不使用的数据)访问，可以认为LDS的容量要比同级的缓存小得多。另外，减少控制逻辑和结构体标记就可以节省出很多LDS空间。

GPU核中允许硬件控制多线程，可以很好的掩盖内存访问延迟。为了让性能和利用率达到最高，就要有足够多的波前执行起来。每个SIMD单元上能运行4个以上(包括4个)波前，或是很多应用需要每个计算单元上能运行16个波前。如果每个SIMD单元可以运行10个波前，那么就有显卡中就会同时有40个活动的波前。为了能更加快速的进行切换，波前的状态则保存在寄存器中，而非缓存中。虽然大量的波前可以掩盖内存访问延迟，但每个波前也都是需要资源支持的，这里必须要进行权衡。

为了减少每个波前产生的请求，系统中的缓存会通过一种过滤机制进行合并读取操作，将写入操作合并，尽可能一次性写入更多的内容——这种机制被称为\emph{合并访问}(coalesing)。向量从规划好的内存(基于DRAM)中读取数据会更加的高效。

图8.6中CPU和GPU是通过PCIe总线连接在一起，两设备间的数据传输则由该总线完成。应为PCIe总线的带宽要比DRAM低很多，和片上内存的带宽没法进行比较。CPU和GPU交互比较多的应用中，带宽会成为应用的性能瓶颈。OpenCL应用中必须最小化与执行相关的数据拷贝。使用离散GPU时，如果应用总要将GPU执行的数据通过PCIe进行拷贝，那么应用就不可能达到最佳的性能。

\subsection{指令集和执行单元}

图8.7展示了基于AMD下一代核芯显示架构Radeon R9 290X中的计算单元。计算单元在该显卡中具有四个SIMD单元。当创建一个波前，其就会被指定到计算单元中的一个SIMD单元上(分配寄存器)，并且在LDS上分配对应尺寸的内存。当波前有64个工作项，SIMD单元具有16个通道，那么SIMD单元上执行一个波前至少需要4个周期。这几个周期都需要一个SIMD单元上完成。这样的话，四个周期之后才能执行新的指令，确切的时间要从前一个指令完全提交到流水线上开始算起。

\myGraphic{0.85}{images/chapter8/8-7.png}{图8.7 Radeon R9 290X计算单元微架构。}

回到计算单元，Radeon R9 290X除了SIMD单元，还有一定数量的执行单元。为了能够进行指令级别的并行，Radeon R9 290X计算单元会查看本轮未向SIMD发射指令的波前，考虑让它们向其他硬件单元发射指令。调度器可能会将5个指令的每个周期放在不同单元上完成，这些单元包括SIMD单元、标量单元、内存单元、LDS或其他指定的功能硬件设备[2]。其他单元的调度与SIMD单元有些不同。比如，标量单元可以完成波前中的一条指令，其设计在每个周期可获得一个新指令。

之前提到的设备，比如第2章提到的Radeon HD 6970，其控制流就由分支单元全权管理。这种比较特殊的执行方式和市面上其他的向量结构有些不同。Radeon R9 290X的设计集成了标量和矢量，就像x86架构的CPU兼容SSE或AVX操作一样。

再回到由64个工作项组成的波前中，每个波前都具有一个指令流(64个工作项共用一个程序计数器)，并且所有分支都以波前为粒度进行执行。为了在Radeon R9 290X上支持分支控制流，架构提供给了执行掩码，用以表示哪些部分写入或不写入结果。因此，波前中工作项间的任何分支，都需要使用\emph{架构指令集}(instruction set architecture，ISA)对指令加以限制，将对应的操作的写入功能开启或关闭。这样执行单元必须指令块队列中的所有可能的分支覆盖到。这样的分支执行会让向量单元执行的效率大大下降。但支持这类控制流免去了程序员手动向量化代码的需要。相同的问题也会出现在其他具有计算能力的架构中，比如NVIDIA的GeForce GTX 780，其软件产品支持宽向量架构，同样也会遇到分支问题。应对分支问题的方式：编程者手动修改代码，编译器对代码进行调整，或是硬件向量化等。

下面的例子会运行在Radeon R9 290X的计算单元上(可见Radeon R9 290X海岛系列ISA标准[3])。下面运行一个简单的内核，在一维索引为0和不为0时所要执行的代码不同：

\begin{codebox}[cpp]
kernel void foo(
  const global int *in,
  global int *out){

  if (get_global_id(0) == 0){
    out[get_global_id(0)] = in[get_global_id(0)];
  } else {
    out[get_global_id(0)] = 0;
  }
}
\end{codebox}

这个简单的内核代码需要使用编译器进行编译，编译出的指令要与ISA进行映射。下面就让我们来看看，编译器是如何在Radeon R9 290X上翻译这段代码的。

\begin{codebox}[text]
s_buffer_load_dword s0, s[4:7], 0x04
s_buffer_load_dword s1, s[4:7], 0x18
s_waitcnt lgkmcnt(0)
s_min_u32 s0, s0, 0x0000ffff
v_mov_b32 v1, s0
v_mul_i32_i24 v1, s12, v1
v_add_i32 v0, vcc, v0, v1
v_add_i32 v0, vcc, s1, v0
s_buffer_load_dword s0, s[8:11], 0x00
s_buffer_load_dword s1, s[8:11], 0x04
v_cmp_eq_i32 s[4:5], v0, 0
s_and_saveexec_b64 s[4:5], s[4:5]
v_lshlrev_b32 v1, 2, v0
s_cbranch_execz label_0016
s_waitcnt lgkmcnt(0)
v_add_i32 v1, vcc, s0, v1
s_load_dwordx4 s[8:11], s[2:3], 0x50
s_waitcnt lgkmcnt(0)
tbuffer_load_format_x v1, v1, s[8:11], 0 offen
    format:[BUF_DATA_FORMAT_32, BUF_NUM_FORMAT_FLOAT]
label_0016:
s_andn2_b64 exec, s[4:5], exec
v_mov_b32 v1, 0
s_mov_b64 exec, s[4:5]
v_lshlrev_b32 v0, 2, v0
s_waitcnt lgkmcnt(0)
v_add_i32 v0, vcc, s1, v0
s_load_dwordx4 s[0:3], s[2:3], 0x58
s_waitcnt vmcnt(0) & lgkmcnt(0)
tbuffer_store_format_x v1, v0, s[0:3], 0, offen
    format:[BUF_DATA_FORMAT_32, BUF_NUM_FORMAT_FLOAT]
s_endpgm
\end{codebox}

清单8.1 Radeon R9 290X上分支操作所生成的ISA指令

这段代码表示在一个通道上所执行的所有操作(一个工作项)。不过，不同于高级语言，这里的标量操作(前缀为\texttt{s\_})会在图8.7所示的GPU标量单元上完成，矢量操作(前缀为\texttt{v\_})会在GPU矢量单元上完成。

清单中第11行(\texttt{v\_cmp\_eq\_i32 s[4:5], v0, 0})，我们能看到这里执行的是矢量比较操作，将工作项ID(存与v0)与常量数0进行比较。比较的结果存于一个64位布尔掩码中，64位掩码由连续的两个标量寄存器组成(\texttt{s[4:5]})。得到的布尔掩码则用来决定OpenCL C内核所要执行的分支。

第12行(\texttt{s\_and\_saveexec\_b64 s[4:5], s[4:5]})，这里隐式的将当前执行掩码与比较结果后的掩码进行与操作。例子中最后得到的执行掩码将用于判断执行if的哪个分支。目标标量寄存器中会将上次执行的掩码进行保存(本例中，\texttt{s[4:5]}将保存上次比较的结果掩码，目前存储的是上次的执行掩码)。当决定哪个工作项走else分支时，就需要使用上次得到的执行掩码。同样，在条件执行完成后需要上次的掩码用来重置执行掩码。这些操作保证对\emph{标量条件代码}(SCC, scalar condition code)寄存器的设置，该寄存器用于触发条件分支。

设置SCC寄存器对程序来说是一种优化，因为其能保证对应工作项进入对应分支，不需要走的分支则不会去执行(本例中，当掩码为0，14行的s\_\allowbreak{}cbranch\_\allowbreak{}execz指令将保证if分支在执行过程中不会触发)。如果条件不满足，代码将进入else分支。本例中，工作项0将使if分支条件成立，其分支掩码为1，所以工作项0则会执行if分支中的语句——用一条SIMD通道完成(从第16行开始(\texttt{v\_add\_i32 v1, vcc, s0, v1}))。

当if分支执行完成，则会在第19行执行一个矢量加载(从tbuffer或纹理内存中加载数据，)，将期望的数据存储到向量寄存器中，v1。第19行是if分支中最后一个操作。注意，在OpenCL C代码中在加载操作完成后，还有一个存储操作。因为，不同分支所要写入的地址是相同的，所以这里编译器将这个存储操作放在程序的最后进行(\texttt{tbuffer\_store\_format\_x v1, v0, s[0:3], 0, offen})。

第22行(\texttt{s\_andn2\_b64 exec, s[4:5], exec})会把原始执行掩码与当前执行掩码的按位取反结果相与。新的执行掩码则表示当前工作项执行的是else分支的语句。else分支不需要从内存中加载数据，工作项只需要将常量0存储v1即可(第23行(\texttt{v\_mov\_b32 v1, 0}))。注意不同的分支都将结果保存在v1寄存器中，这样就可以将存储操作放在之后完成。

细心的读者可能已经发现了，指令中这里并没有对else分支进行忽略。本例中，编译器认为与其忽略else分支，不如不对其进行加载，这样就节省了执行预测的性能开销。而且，这样操作就不会对v1寄存器进行更新，和忽略else分支达到了同样的效果。

这是一个简单的例子。在当有嵌套或深层条件时，就会有更加复杂和更加长的序列用来存储条件掩码，并且与执行掩码进行与操作，逐级缩小执行管道最后所要执行的变量分支。逐级缩小的过程将会对造成执行性能的下降。所以，良好的代码设计将会在执行相同指令的前提下，获得更佳的性能。程序的掩码管理十分复杂，并且GPU向量操作与CPU向量操作(SSE)也区别。使该ISA区别于SSE等CPU ISA的是精细的掩码管理与向量操作，而非核数更多。

最后，第24行(\texttt{s\_mov\_b64 exec, s[4:5]})，将执行掩码进行重置，并且将数据存储在v1寄存器中，在第30行(\texttt{tbuffer\_store\_format\_x v1, v0, s[0:3], 0, offen})使用tbuffer将数据存储到对应位置上。

\subsection{资源分配}

GPU中的每个SIMD单元都有固定数量的寄存器和LDS存储空间。每个计算单元中有256KB的空间供寄存器使用。这些寄存器分为四组，每个SIMD单元有256个寄存器，每个寄存器，每个单元有64个通道，每个通道具有32位。计算单元中，寄存器按计算单元上运行的波前数量进行划分。每个计算单元有64KB的LDS，就像对32行的SRAM一样可以进行随机访问。LDS会在计算单元中分成多份，划分给对应的工作组，用于开辟对应的局部内存(通过OpenCL运行时传入的参数进行分配)。

计算单元执行内核时，我们可能看到如图8.8所示的资源瓶颈。图中两个工作组，每个工作组包含两个波前，每个工作项需要42个向量寄存器(按照波前的尺寸等比例放大)，并且共享50个标量寄存器，并且每个工作组需要24KB LDS。每个计算单元能够同时执行4个这样的波前，这样可将每个波前分配给一个SIMD单元，这一数量受工作组LDS需求限制，低于让设备保持忙碌所需的最小波前数。不过，当执行标量代码或内存操作时，就没有空间进行切换操作。下面给出在Radeon R9 290X上工作组对一个计算单元的占用比(等式8.1)：

\myFormula{0.62}{images/chapter8/8-1-exp.png}{8.1}

\myGraphic{0.85}{images/chapter8/8-8.png}{图8.8 将OpenCL内存模型映射到Radeon R9 290X GPU。}

等式8.1中，VGRP表示矢量通用寄存器，SGPR表示标量通用寄存器，WI表示工作项，WG表示工作组，WF表示波前。通过OpenCL内核参数传入的部分使用粗体表示。

图8.8的例子中，如果增加SIMD单元中波前运行的数量(4个或更多)，在控制流和内存延迟方面，让标量和矢量单元忙碌起来，有助于对延迟的隐藏。本例中LDS有限，对于算法来说，将波前数量(只有一个波前)增加到3个是个不错的选择。另外，减少LDS内存分配，则可以在每个计算单元上再运行第三个工作组。这个方式在波前等待栅栏或内存访问时非常有用(SIMD单元闲置状态)。

每个波前运行在一个SIMD单元上，直到执行完成。任意组具有相同工作组的波前都会在同一个计算单元上完成。由于每个工作组的状态存储量的原因，才将这些工作放在同一个计算单元完成。本例中，我们看到每个工作组需要24KB LDS，且寄存器可操作的空间超过21KB。在刷新内存和搬移到其他核上时，这个数据量还是相当大的。当内存控制器执行高延迟的读或写操作时，如果没有其他波前在计算逻辑单元(ALU)中时，那就没有任务分配到SIMD单元上，计算硬件这时处于闲置状态。

\myRefRule

\begin{myRefs}

[2] Advanced Micro Devices, GCN Whitepaper. AMD Graphics Core Next(GCN) Architecture, 2013. \texthref{http://www.amd.com/Documents/GCN\_Architecture\_whitepaper.pdf}

[3] Advanced Micro Devices, ADM Sea Islands Series Instruction Set Architecture, 2013. \texthref{http://developer.amd.com/wordpress/media/2013/07/AMD\_Sea\_Islands\_Instruction\_Set\_Architecture.pdf.}

\end{myRefs}



% ---- 8.3-chinese.md ----
\section{OpenCL内存性能的考量}

\subsection{全局内存}

与内存相关的话题我们在第2章有详细的讨论过。OpenCL应用的性能与是否高效的使用的内存有着很大的关系。不过，高效的内存则依赖与具体的硬件(执行OpenCL内核的设备)。因此，同样的访存模式，在GPU上是高效的，不过在CPU上就不一定了。因为GPU的供应商繁多，且与CPU在制造方面有很大的区别。

所有例子中，内核吞吐量的级别是内存性能分析的开端。下面简单的计算公式，就是用来计算内核的带宽大小：

\myFormula{0.62}{images/chapter8/8-2-exp.png}{8.2}

EB代表有效带宽，Br代表从全局内存上读取的数据量(单位：byte)，Bw表示写入全局内存的数据量(单位：byte)，t代表内核运行的时间。

时间t的获取，可以通过一些性能测评工具，比如：ADM的CodeXL。Br和Bw可以通过每个工作项所读取或写入的数据量，然后乘以工作项的数量计算得出。所以，在某些情况下，这些读写的数据量都是估算出来的。

当我们测得当前执行内核的带宽，我们可以将测出带宽与执行设备的峰值带宽进行比较，看一下二者在数值上差多少。如果二者很接近，那说明我们充分的利用了当前内存系统；如果二者相差甚远，那我们就需要考虑重构我们内存访问的方式，以提高内存的利用率，增大内核带宽。

空间局部性是OpenCL访存中需要重点考虑的问题。大多数架构在运行OpenCL内核时，会基于不同等级的矢量进行(可能像SSE，或使用管道导向型输入语言进行在佛那个向量化，例如AMD的IL或NVIDIA的PTX)，内存系统会将向量中的数据打包一起处理以加速应用。另外，局部访问通常会使用缓存进行。

大多数现代CPU都支持不同版本的SSE和AVX向量指令集。将内存部分设计成全对齐的模式，向量读取这样的内存会有使用到相关的指令集，并且向量指令会使内存访问更加高效。可以给定一个较小的向量尺寸(比如float4)，这样编译器会生成更加高效的向量读取指令。其很好的利用了缓存行，缓存和寄存器间做数据移动是最高效的。不过，CPU在处理未对齐的内存或更多的随机访问时，缓存会帮助掩盖一些性能损失。图8.9和图8.10提供了两个例子，一个用于读取一段连续的4个数据，另一个则是通过随机访问4个数据。如果缓存行较窄，则会出现很多次缓存未命中，这样的情况会大大影响应用的性能。

\myGraphic{0.85}{images/chapter8/8-9.png}{图8.9 使用向量读取能更高效地让数据从内存系统返回；当工作项访问连续元素时，GPU硬件通过合并访问也能达到同样的效果。}

\myGraphic{0.85}{images/chapter8/8-10.png}{图8.10 访问非连续元素时只能返回较小的数据块，效率更低。}

之前的章节中我们曾讨论过，GPU内存架构与CPU的内存架构有着很大的不同。GPU使用多线程的方式来掩盖不同级别的内存延迟，GPU更偏向于ALU能力，而不是缓存和复杂的乱序逻辑。传统GPU具有更多的计算资源可用，如果我们不想GPU饿死，就需要具有更加高带宽的内存系统。很多现代GPU架构，特别是高性能桌面版本，比如AMD的Radeon系列和NVIDIA的GeForce系列，都在使用宽SIMD架构。试想8.10图中的例子，将向量扩展成(AMD Radeon R9支持的)64位硬件向量。

高效的访问在不同的架构中也有不同的方式。对于x86 CPU来说使用的是SSE指令集，我们可能会使用128位的float4类型作为数据处理的单元，这样可能会增加缓存行的利用率，减少缓存未命中的概率。对于AMD Radeon R9 290X GPU架构，同一波前中连续的工作项就可以同时对内存进行访问。如果内存系统不能及时的处理这些请求，则会造成访存延迟。对于最高的性能，同一波前的工作项同时发起32位的读取请求时，就意味着最多需要读取256字节(32bit x 64个工作项)内存数据，这样内存系统只需要相应开辟一块较大内存上的请求。为了可在不同的架构间进行移植，一个好的解决方式就是让内存访问的效率尽可能的高，可以在宽矢量设备(AMD和NVIDIA GPU)和窄矢量设备(x86 CPU)都有很好的访存效率。为了达到这种效果，我们可以通过工作组计算出所要放访问的内存起始地址，该地址应为work-groupSIze x loadSize对齐，其中loadSize是每个工作项所加载的数据大小，其应该是一个合理的值——对于AMD GCN架构的设备来说，32位是一个不错大小；对于x86 CPU和旧一些的GPU架构来说，128位则是很好的选择；对支持AVX架构的设备来说，256位则是不二之选。为什么32位对于AMD GCN架构来说是一个不错的选择，下面我们就来解释一下。

要处理不同的内存系统，需要面对很多问题，比如：减少片外链接DRAM的访存冲突。先让我们来看一下AMD Radeon架构中如何进行地质分配。图8.11中低8位展示了给定内存块中的数据内存起始地址；这段信息可以存储在缓存行和子缓存行中，供我们进行局部读取。如果尝试读取二维数据中一列的数据，对于行优先的存储方式，这种方式对于片上总线来说是低效的。其也意味着，设备上同时执行的多个工作组，访问的内存通道和内存块都有所不同。

\myGraphic{0.85}{images/chapter8/8-11.png}{图8.11 把Radeon R9 290X的地址空间映射到内存通道和DRAM块。}

\myGraphic{0.85}{images/chapter8/8-12.png}{图8.12 Radeon R9 290X内存子系统。}

每个内存通道的控制器与片外内存进行连接(图8.12)。我们希望执行设备能够通过内存系统，访问到其中的所有内存块与内存通道。不过，一个波前中某个矢量将会命中多个内存通道(或内存块)占用和阻塞他们，导致其他波前对相应地址访问的延迟，从而导致带宽的下降。最佳的方式是，给定的波前能够连续的读取指定内存通道和内存块上的数据，允许更多波前可以并行的对内存进行访问，形成高效的数据流。

为了避免使用多个通道，一个波前所能访问到的区域在64个字(256字节)之内，这就能保证所有工作项在读取32位数据时，是从连续的地址上获取。这里需要讨论一下访存最坏的结果——"当多个波前中的每个工作项对同一地址的数据进行访问"——每个工作项中的变量都将命中同一内存通道和内存块，并且串行的获取数据，访存吞吐量比顶峰时降低数倍。更多的有关AMD架构主题的内容，可以在《AMD的OpenCL编程指南》中找到[4]。同样的信息也在其他GPU供应商的计算架构中出现——例如，NVIDIA的《CUDA编程指南》[5]。

\subsection{局部内存——软件可控缓存}

大多数支持OpenCL的设备都支持不同形式的缓存。由于面向图像的设计，很多GPU只提供只读数据缓存，这样能复用大量的数据。

OpenCL使用缓存最简单的方式就是使用图像类型(第6和第7章讨论过)。GPU上，图像可以将数据映射成硬件可读纹理。假设复杂的滤波器不需要进行二维内存访问，那么访问内存效率将会提高。不过，GPU缓存相较波前所要读取的内存，则是相形见绌。编程者可通过代码控制暂存式内存在局部空间的大小，在高效获取数据的同时，减少了硬件控制缓存的开销，有效的节约资源。这对于工作组内工作项的数据交换来说，能够减少栅栏冲突和访问延迟。(图5.5就是一个例子)

当然，考虑对数据进行优化时，需要认真的考虑如何利用数据的局部性。很多例子中，其消耗在于使用额外的拷贝指令将数据搬移到局部内存中，之后搬运到ALU中(可能通过寄存器)进行计算，这种方式的效率通常要比简单的复用缓存中数据的效率低得多。大量的读取和写入操作复用同一地址时，将数据搬移到局部内存中将会很有用，读取和写入操作对局部内存的操作延迟要远远小于对于全局内存的操作。并且，对二维数据进行访问时，就不需要通过全局变量进行数据加载，从而减少缓存加载所需要的时间。

下例中的读/写操作，将会大大收益与局部内存，特别是给定宽度的只读缓存。读者可以尝试，写出前缀求和的C代码和下面的代码进行对照：

\begin{codebox}[cpp]
void localPrefixSum(
  __global unsigned *input,
  __global unsigned *output,
  __local  unsigned *prefixSums,
           unsigned numElements){

  /* Copy data from global memory to local memory */
  for (unsigned index = get_local_id(0);
    index < numElements;
    index += get_local_size(0)){
    prefixSums[index] = input[index];
  }

  /* Run through levels of tree, each time halving the size
   * of the element set performing reduction phase */
  int offset = 1;
  for (unsigned level = numElements / 2;
    level > 0;
    level /= 2){
    barrier(CLK_LOCAL_MEM_FENCE);

    for (int sumElement = get_local_id(0);
        sumElement < level;
        sumElement += get_local_size(0)){
        int ai = offset * (2 * sumElement + 1) - 1;
        int bi = offset * (2 * sumElement + 2) - 1;
        prefixSums[bi] = prefixSums[ai] + prefixSums[bi];
    }
    offset *= 2;
  }

  barrier(CLK_LOCAL_MEM_FENCE);

  /* Need to clear the last element */
  if (get_local_id(0) == 0){
    prefixSums[numElements - 1] = 0;
  }

  /* Push values back down the tree */
  for (int level = 1; level < numElements; level *= 2){
    offset /= 2;
    barrier(CLK_LOCAL_MEM_FENCE);

    for (int sumElement = get_local_id(0);
        sumElement < level;
        sumElement += get_local_size(0)){
        int ai = offset * (2 * sumElement + 1) - 1;
        int bi = offset * (2 * sumElement + 2) - 1;
        unsigned temporary = prefixSums[ai];
        prefixSums[ai] = prefixSums[bi];
        prefixSums[bi] = temporary + prefixSums[bi];
    }
  }

  barrier(CLK_LOCAL_MEM_FENCE);

  /* Write the data out to global memory */
  for (unsigned index = get_local_id(0);
    index < numElements;
    index += get_local_size(0)){
    output[index] = prefixSums[index];
  }
}
\end{codebox}

代码清单8.2 单工作组的前缀求和

尽管这段代码对许多架构来说并非最优，其让工作组中的工作项共享一段数组，从而降低访问数据的延迟。数据流中的第一个循环(第19行)如图8.13所示。注意循环中的每次迭代更新一些工作项中的数据，以供下次迭代使用。这里需要工作项之间的合作，才能完成这项工作。内部的循环最好能覆盖大部分工作项，以避免执行分支。为了保证工作项的正确行为，我们在外层循环中添加了栅栏，同步对应的工作项，以确保数据在下次循环时已经准备好。

\myGraphic{0.85}{images/chapter8/8-13.png}{图8.13 代码清单8.2中前缀和的累加过程：16个元素放在局部内存中，由8个工作项完成。}

代码清单8.2中的前缀求和使用的是局部内存的方式，这种方式在大多数宽SIMD架构下的效率并不高(比如：高端GPU)。之前我们讨论过全局变量，内存系统倾向于将内存进行分块，而非让每个内存位置都能让外部访问到。最后，硬件的暂存式内存(比如：缓存)更加倾向于让每个内存块执行多次读操作，或者并发的进行读写操作(或其他一些多次访问操作)，多次读取操作可以跨越多个内存块进行。所以在使用宽硬件SIMD访问内存时，这点必须要考虑到。每个周期中，Radeon R9 290X GPU使用4个SIMD单元可以处理两个局部内存操作。每个SIMD单元具有16个通道，每个周期可以完成32个读或写操作，可以LDS上32个内存块。每个内存块支持一个访问入，那么每个内存块只能提供一个值，所以只有在所有访问目标都在不同的内存块上时，才能获取最好的吞吐率。同样的规则也适用于其他计算架构：NVIDIA的费米框架，局部内存分为32个内存块。

局部内存的问题并没有全局内存那么严重。全局内存中，大范围跨越访问会因缓存行上的多个访问丢失，导致访存延迟。局部内存中，至少架构中是由暂存器的，编程者可以根据自身意愿将数据放到该内存上。这里需要的是，我们发起的16次局部内存访问，最好访问的是不同的内存块。

图8.14中展示了图8.13前缀求和在8个LDS块上进行的第一步，这里每个工作项在每个周期中能执行一次局部内存操作。这个例子中，我们的局部内存可以在一个周期内返回8个值。那么我们应该如何在第一步中提升应用的执行效率呢？

注意16个元素的局部内存(对于前缀求和来说是很有必要的)已经超过两行。每列上的数据在一个内存块上，每行上只存放每个内存块的一个地址。假设(通常在很多架构中)每个内存块都是32位宽，并且假设当前波前不需要依赖其他波SIMD单元处理后的数据。不过在图8.14中，我们的第一步可谓是失败的。局部内存和全局内存一样，一个SIMD矢量会根据矢量长度连续，并高效的访问数据，并且局部内存不存在访问争夺。图8.13中我们看到的是另一番景象。29行中prefixSums[bi]=perfixSums[ai]+perfixSums[bi]中，对prefixSums[bi]访问了两次。图中就是尝试对位置3、7、11和15进行读取。图8.14中，3和11都存储在内存块3上，7和15存储在内存块7上。根据之前的描述，就不要想同时读取到内存块上的两个值，所以访问同一块内存的操作只能在GPU上串行执行，从而导致读取延迟。为了达到最优性能，我们要尽可能避免冲突。有个十分有用的技巧是进行简单的填充地址，就如图8.15所示。为了偏移地址(以内存块大小对齐)，甚至可以改内存访问的跨度。不过，对地址操作的开销要比内存块冲突严重的大多；因此，我们要在具体的架构上进行调试。

\myGraphic{0.85}{images/chapter8/8-14.png}{图8.14 图8.13的第一步：具有8个存储体的LDS的工作情况。}

\myGraphic{0.85}{images/chapter8/8-15.png}{图8.15 图8.14的第一步：对原始数据做填充(padding)以消除LDS中的存储体冲突。}

局部内存的大小是固定的。任何OpenCL设备都具有一段暂存式内存，这块内存大小有限，且不归硬件管理。Radeon R9 290X GPU式例中，其局部空间大小为64KB。要注意的是这64KB是所有工作组一共可以使用的共享内存。因为GPU是用多线程提高吞吐量，从而掩盖访存的延迟。如果每个工作组使用16KB，那么每个核上只能运行4个工作组。如果几个波前(一个或两个)中包含有几个工作组，那么这样就刚好能够掩盖访问延迟。虽然局部内存能提升应用性能，不过局部内存大小有限。因此，我们需要在使用局部内存和减少硬件线程方面进行平衡。

OpenCL运行时API也支持查询对应设备上局部内存的大小。在编程者编译或暂存局部内存数据时，其可以作为OpenCL内核参数。下面代码第一个调用，是用来查询局部内存类型，其可以用来判断哪些内存属于局部内存或全局内存(哪些是可以缓存或不能缓存的)，第二个调用时用来返回局部内存的大小：

\begin{codebox}[cpp]
cl_int err;
cl_device_local_mem_type type;
err = clGetDeviceInfo(
  deviceId,
  CL_DEVICE_LOCAL_MEM_TYPE,
  sizeof(cl_device_local_mem_type),
  &type,
  0);

cl_ulong size;
err = clGetDeviceInfo(
  deviceId,
  CL_DEVICE_LOCAL_MEM_SIZE,
  sizeof(cl_ulong),
  &size,
  0);
\end{codebox}

\myRefRule

\begin{myRefs}

[4] Advanced Micro Device, The AMD Accelerated Paralel Processing-OpenCL Programming Guide, Advanced Micro Devices, Inc,. Sunnyvale, CA, 2012.

[5] NVIDIA, CUDA C Programming Guide, NVIDIA Corporation, Santa Clara, CA, 2012.

\end{myRefs}



% ---- 8.4-chinese.md ----
\section{本章总结}

本章旨在表示OpenCL在具体架构上的映射关系。本章中，我们简单的比较了一下CPU和GPU架构的不同，不同的矢量宽度带来的巨大的性能差异(NVIDIA的GPU矢量宽度为32位，AMD的GPU矢量宽度是64位，CPU则小的多)，还有不同的线程管理和指令调度。本书中我们不能覆盖所有市面上的架构，不过本章中我们给出了一些例子。所以，编程者在以某个OpenCL设备作为目标设备时，就需要仔细研究对应平台架构的相关文档，这样才能在对应平台上获得最优的性能。


